% This is samplepaper.tex, a sample chapter demonstrating the
% LLNCS macro package for Springer Computer Science proceedings;
% Version 2.21 of 2022/01/12
%
\documentclass[runningheads]{llncs}
%
\usepackage[T1]{fontenc}
% T1 fonts will be used to generate the final print and online PDFs,
% so please use T1 fonts in your manuscript whenever possible.
% Other font encondings may result in incorrect characters.
%
\usepackage{graphicx}
% Used for displaying a sample figure. If possible, figure files should
% be included in EPS format.
%
% If you use the hyperref package, please uncomment the following two lines
% to display URLs in blue roman font according to Springer's eBook style:
%\usepackage{color}
%\renewcommand\UrlFont{\color{blue}\rmfamily}
%\urlstyle{rm}
%

\usepackage{booktabs}
\usepackage{amsmath}
\usepackage{hyperref}

\usepackage{subfigure}

\usepackage[table,xcdraw]{xcolor}
\usepackage{colortbl}

\usepackage{flushend}

\usepackage{listings}
\usepackage{xcolor}
% \usepackage[hidelinks]{hyperref}
% \lstset{
%   basicstyle=\ttfamily\footnotesize,
%   keywordstyle=\color{blue},
%   commentstyle=\color{gray},
%   numbers=left,
%   numberstyle=\tiny,
%   numbersep=5pt,
%   frame=single,
%   columns=flexible,
%   keepspaces=true,
%   showstringspaces=false,
%   captionpos=b,
%   escapeinside={(*@}{@*)}
% }

% Para notas TODO (remover na versão final)
\newcommand{\todo}[1]{\textcolor{red}{[TODO: #1]}}
\newcommand\ms[1]{{\color{blue}\textbf{?MS: }#1}}
\newcommand\rem[1]{{\color{red}\textbf{?REMOVE: }#1}}
\begin{document}
%
%==============================================================================
% TITLE OPTIONS (choose one)
%==============================================================================
% Option A: Focus on gap characterization
% \title{Mind the Gap: Characterizing GPU Data Compression\\Performance Across AMD Architectures}
\title{Characterizing GPU Data Compression Performance Across AMD Architectures}

% Option B: Focus on porting
%\title{From CUDA to HIP: Porting nvCOMP Compression\\Algorithms to AMD GPUs}

% Option C: Focus on cross-generation
%\title{Cross-Generation Analysis of LZ4 Compression\\on AMD CDNA and RDNA GPUs}

% Option D: Technical focus
%\title{Bridging the Compression Gap: An Evaluation of\\Lossless Data Compression on AMD GPU Platforms}

%
% \titlerunning{Abbreviated paper title}
% If the paper title is too long for the running head, you can set
% an abbreviated paper title here
%
\author{Cristiano Künas\inst{1}\orcidID{0000-0003-1080-7230} \and
Gabriel Freytag\inst{1}\orcidID{0000-0001-6081-5904} \and
% Alexandre Sardinha\inst{2}\orcidID{0009-0003-0418-3289} \and
% Alexandre Carissimi\inst{1}\orcidID{0000-0002-0884-1483} \and
Philippe Navaux\inst{1}\orcidID{0000-0002-9957-5861} 
}
%
\authorrunning{C. Künas et al.}
% First names are abbreviated in the running head.
% If there are more than two authors, 'et al.' is used.
%
\institute{Federal University of Rio Grande do Sul (UFRGS), Porto Alegre, Brazil \\
\email{\{cakunas, gfreytag, navaux\}@inf.ufrgs.br}
% \and 
% Petroleo Brasileiro S.A., Rio de Janeiro, Brazil \email{alexandre.sardinha@petrobras.com.br} 
% \and
% ABC Institute, Rupert-Karls-University Heidelberg, Heidelberg, Germany\\
% \email{\{abc,lncs\}@uni-heidelberg.de}
}
%
\maketitle              % typeset the header of the contribution
%

%==============================================================================
\begin{abstract}
GPU-accelerated data compression is essential for mitigating I/O bottlenecks in high-performance computing, yet a significant tool gap exists between NVIDIA and AMD platforms. While NVIDIA's nvCOMP library provides highly optimized compression algorithms, AMD users lack equivalent, mature solutions — a pressing concern given AMD's growing presence in top HPC systems such as Frontier and LUMI. This paper addresses this gap through two contributions. First, we present a quantitative characterization of LZ4, Snappy, and Cascaded compression performance across three generations of AMD GPUs (MI50, MI210, MI300X) and the RDNA~3-based RX~7900~XT. Second, we introduce an initial port of the nvCOMP~2.2 algorithms to HIP, enabling compression on AMD platforms. Our evaluation using synthetic and scientific datasets reveals that: (1)~the MI300X achieves up to 11$\times$ higher decompression throughput than the MI50; (2)~RDNA~3 surprisingly outperforms CDNA~3 for compression workloads; and (3)~while a performance gap relative to optimized NVIDIA solutions remains, GPU compression on AMD is already beneficial for I/O-bound applications. Our results establish baseline metrics for AMD GPU compression and identify optimization opportunities for future work.

\keywords{GPU Computing \and Data Compression \and LZ4 \and AMD \and HIP \and ROCm \and High-Performance Computing}
\end{abstract}

%
%
%


%\section{Introduction}
\label{sec:introduction}

Modern scientific computing faces an increasingly critical
challenge: the volume of data generated by simulations is growing
faster than the ability to move it efficiently. High-performance
computing applications such as seismic forward modelling, climate
simulation, and molecular dynamics produce datasets ranging from
gigabytes to petabytes, making data movement the dominant
bottleneck in many workflows~\cite{jin2024concealing,lee2002migration}.
GPU-accelerated compression has emerged as an attractive solution,
leveraging the parallel architecture of graphics processors to
achieve throughputs that can exceed storage system
bandwidths~\cite{nvcomp,knorr2021ndzip}.

The high-performance computing landscape, however, has become
increasingly heterogeneous. AMD GPUs now power three of the top ten
systems on the TOP500 list, including the two highest-ranked
exascale systems (El Capitan and Frontier)~\cite{top500nov2025}.
Despite AMD's growing presence in HPC, the availability of optimized
GPU compression tools has not kept pace with the hardware. The
NVIDIA ecosystem benefits from \nvcomp~\cite{nvcomp}, a
vendor-supported, production-ready library available since 2020
that offers architecture-tuned implementations of LZ4, Snappy,
Cascaded, GDeflate, ANS, BitComp, and (since release 5.x) ZFP. In
contrast, the HIP/ROCm ecosystem relies on community-driven efforts
such as \hipcomp{}, which appeared in 2022 and remains experimental,
with limited algorithm coverage, direct-port optimization, and
sparse documentation. This asymmetry creates a tooling vacuum for
the growing number of HPC systems built on AMD GPUs.

A functional hipify-based port of \nvcomp{} establishes an
AMD-compatible baseline~\cite{anon_iccsa2026,kunas2025compressao},
but the resulting implementation is far from optimal: it inherits
launch geometries and host-to-device staging strategies that were
tuned for NVIDIA hardware, leaving substantial throughput on the
table. Recent work on seismic reverse-time migration pipelines
also showed that, even with a correctly ported library, the I/O
subsystem dominates end-to-end time~\cite{kunas2024carla,
boito2017seismic}, which makes the host-to-device staging on the
input path a high-leverage optimization target. Two observations
from a systematic cross-architecture characterization campaign on
AMD's CDNA and RDNA lineups motivate the present work. First, at the kernel level, the default chunk size of 64~KB
inherited from \nvcomp{} leaves the larger AMD parts (notably the
MI300X) wave-count-starved: the number of chunks dispatched is one
to two orders of magnitude smaller than the GPU's wave slot
capacity, and the SIMT pipelines run nearly idle. Second, at the
transfer level, the dominant cost in end-to-end execution is the
host-to-device staging of the input batch, which a naive
per-chunk \texttt{hipMemcpyAsync} from pageable storage caps at
roughly 6~GB/s on PCIe Gen4 -- a fifth of the bus theoretical
peak.

This paper presents \arcto, a profile-driven GPU compression
library for AMD that closes both gaps with two complementary
contributions:

\begin{itemize}
\item A \textbf{wave-saturation chunk-size formula}, derived from
hardware performance counter analysis, that picks the chunk size
which keeps every wave slot busy and unlocks 1.4--2.9x additional
throughput on the compress kernels with no code change other than
a single launch-time parameter.

\item An \textbf{adaptive tiled aggregation API}, derived from a
cost model fitted to the characterization data, that replaces the
single-shot coalesced pinned-host staging used by the published
ICCSA baseline. We show that the single-shot path, when measured
honestly with full lifecycle accounting (\texttt{hipHostMalloc}
plus pageable-to-pinned memcpy plus bulk H2D), is in fact slower
than the scattered-pageable baseline at every input size from
100~MB to 4~GB on RDNA3. The adaptive variant amortizes the
allocation cost across windows of a profile-driven optimal size,
recovering the gain and reaching \textbf{up to 2.26x end-to-end
speedup} over the baseline on MI300X for the 16~GB TTI workload.
\end{itemize}

Both optimizations are derived from performance counter observations
on the AMD lineup, not from arch-agnostic intuition, and both
generalize across the four AMD generations evaluated (MI50, MI210,
MI300X, RX~7900~XT). The cost model behind the adaptive contribution
is also validated by direct measurement: the optimal window size
picked on each architecture by the auto-detection rule matches the
empirically best window size within one octave.

The remainder of the paper is organized as follows.
Section~\ref{sec:background} surveys the AMD GPU landscape, the
batched compression model used by \nvcomp{} and inherited by
\arcto, and the experimental setup used throughout the evaluation.
Section~\ref{sec:arcto-library} \TODO{describes the \arcto{} library
design and its hipify-based porting strategy}.
Section~\ref{sec:optim-chunk-size} \TODO{presents the wave-saturation
chunk-size formula}.
Section~\ref{sec:optim-adaptive-tiled} \TODO{presents the adaptive
tiled aggregation contribution with cost model derivation}.
Section~\ref{sec:evaluation} \TODO{reports the cross-architecture
scaling sweep}.
Section~\ref{sec:related-work} \TODO{positions the work in the
literature}, and Section~\ref{sec:conclusion} concludes.


\section{Introduction}
\label{sec:introduction}

Modern scientific computing faces an increasingly critical challenge: the volume of data generated by simulations is growing faster than our ability to move it efficiently. High-performance computing (HPC) applications in seismic imaging, climate modeling, and computational fluid dynamics routinely produce datasets ranging from terabytes to petabytes, making data movement---rather than computation---the dominant bottleneck in many workflows. This I/O wall has motivated significant interest in data compression techniques that can reduce storage requirements and transfer times without sacrificing scientific fidelity.

GPU-accelerated compression has emerged as a particularly attractive solution, leveraging the massively parallel architecture of graphics processors to achieve throughputs that can exceed storage system bandwidths. On NVIDIA platforms, the nvCOMP library~\cite{nvcomp} has become the de facto standard, offering highly optimized implementations of algorithms such as LZ4, Snappy, and Cascaded that achieve compression throughputs exceeding hundreds of GB/s on modern GPUs. This mature ecosystem has enabled researchers to seamlessly integrate high-performance compression into their CUDA-based workflows.

However, the HPC landscape is becoming increasingly heterogeneous. AMD GPUs now power several of the world's most powerful supercomputers, including Frontier at Oak Ridge National Laboratory---currently the world's fastest supercomputer---and LUMI at CSC Finland, which ranks among the top systems in the TOP500 list. This shift creates a pressing need for portable, high-performance software tools that can execute efficiently across different GPU vendors. Unfortunately, a significant \textbf{tool gap} exists: while NVIDIA users benefit from years of optimization in nvCOMP, AMD users lack equivalent mature compression libraries. Existing solutions such as hipCOMP provide basic functionality but have not received the same level of optimization effort, leaving researchers with AMD hardware at a disadvantage.

This paper addresses this gap through two contributions. First, we present a \textbf{quantitative characterization} of compression performance across three generations of AMD GPUs, spanning the CDNA and RDNA architectures. Using both synthetic datasets and scientific data representative of seismic applications, we establish baseline performance metrics and analyze how compression throughput scales with architectural improvements. Second, we introduce an \textbf{initial porting effort} that brings key compression algorithms from nvCOMP 2.2 to AMD platforms via the HIP programming model. While this port prioritizes functional correctness over AMD-specific optimizations, it demonstrates that the performance gap can be mitigated and provides a foundation for future optimization work.

Our experimental evaluation spans four AMD GPUs: the MI50 (CDNA/Vega), MI210 (CDNA~2), MI300X (CDNA~3), and RX~7900~XT (RDNA~3). We evaluate three compression algorithms---LZ4, Snappy, and Cascaded---using datasets that range from highly compressible (zeros) to realistic scientific data (TTI wavefields from seismic simulations). Our results reveal several key findings:

\begin{enumerate}
    \item AMD GPUs have improved dramatically across generations, with the MI300X achieving up to 11$\times$ higher decompression throughput than the MI50 baseline.
    \item Architectural differences between CDNA and RDNA lead to distinct performance characteristics, with RDNA~3 surprisingly outperforming CDNA~3 for certain compression workloads.
    \item While a performance gap relative to nvCOMP remains, compression on AMD GPUs is already beneficial for I/O-bound workloads, and significant optimization opportunities exist.
\end{enumerate}

The remainder of this paper is organized as follows. Section~\ref{sec:background} provides background on GPU-based compression and motivates the need for AMD-compatible tools. Section~\ref{sec:implementation} describes our porting approach and implementation. Section~\ref{sec:methodology} details the experimental methodology. Section~\ref{sec:results} presents our cross-generation performance analysis. Section~\ref{sec:discussion} discusses implications and future directions. Section~\ref{sec:related} reviews related work, and Section~\ref{sec:conclusion} concludes the paper.





%\section{Background and Experimental Setup}
\label{sec:background}

\subsection{AMD GPU Architectures}
\label{sec:bg-amd}

AMD has established two distinct GPU architectures targeting
different markets. The Compute DNA (CDNA) line is purpose-built for
data-center HPC, emphasizing high computational throughput and
high-bandwidth memory; the evolution from CDNA (MI50, MI100)
through CDNA~2 (MI210, MI250X) to CDNA~3 (MI300X) brought
substantial improvements in compute density and HBM bandwidth, with
the current flagship MI300X providing 304 compute units, 192~GB of
HBM3, and over 5~TB/s of memory bandwidth. The Radeon DNA (RDNA)
line targets gaming and consumer applications but is also relevant
for cost-sensitive HPC deployments and for development workflows;
RDNA~3 (RX~7900~XT) is the reference part here. Table~\ref{tab:gpus}
summarizes the four AMD GPUs evaluated in this paper.

\begin{table}[t]
\centering
\caption{AMD GPU architectures evaluated in this paper.}
\label{tab:gpus}
\footnotesize
\begin{tabular}{lllrrr}
\toprule
\textbf{GPU} & \textbf{Arch.} & \textbf{ISA} & \textbf{CUs} & \textbf{Memory} & \textbf{BW (GB/s)} \\
\midrule
MI50       & GCN    & gfx906  &  60 & 32 GB HBM2  & 1{,}024 \\
MI210      & CDNA~2 & gfx90a  & 104 & 64 GB HBM2e & 1{,}638 \\
MI300X     & CDNA~3 & gfx942  & 304 & 192 GB HBM3 & 5{,}300 \\
RX 7900 XT & RDNA~3 & gfx1100 &  84 & 20 GB GDDR6 &   800  \\
\bottomrule
\end{tabular}
\end{table}

The four reference parts span an order of magnitude in compute unit
count and memory bandwidth, and three distinct architectural
families. They share the SIMT execution model but differ on the
specific axes that compression kernels exercise the most: wavefront
size (CDNA fixed at 64 threads; RDNA configurable wave32 or wave64),
on-package cache capacity (4~MB on MI50 to 256~MB Infinity Cache on
MI300X), and host interconnect (PCIe Gen3 to Gen5 across the
lineup). These differences propagate directly into the
optimization choices presented in Sections~\ref{sec:optim-chunk-size}
and~\ref{sec:optim-adaptive-tiled}.

\subsection{Batched Chunked Compression on GPUs}
\label{sec:bg-batched}

The dominant GPU compression model, introduced by \nvcomp{} and
inherited by \arcto, is \emph{batched chunked} processing. The
input buffer is partitioned into fixed-size chunks (typically
64~KB) and the chunks are processed independently in parallel,
with one or more waves per chunk according to the algorithm's
internal structure. LZ4 uses a single wave per chunk with a
private hash table for match finding; Snappy uses two
cooperating waves; Cascaded applies a multi-stage pipeline
(delta plus run-length plus bit-packing) with 128 threads per
4~KB partition. The chunk size is the central parameter exposed
to the caller and is the primary lever for occupancy tuning.

At the API boundary, the caller hands the library a vector of
chunk pointers and an output device buffer. Two transfer
patterns are common: \emph{scattered pageable}, where each chunk
is staged from a separate pageable host allocation via its own
\texttt{hipMemcpyAsync}, and \emph{coalesced pinned}, where the
chunks are first packed into a single pinned host buffer and the
entire batch is uploaded in one bulk \texttt{hipMemcpyAsync}. The
pinned variant achieves PCIe peak bandwidth on the bulk copy but
pays a non-trivial allocation cost up front. The trade-off
between these two patterns is the subject of
Section~\ref{sec:optim-adaptive-tiled}.

\subsection{The Tooling Gap on AMD}
\label{sec:bg-gap}

The AMD software ecosystem centers on ROCm~\cite{rocm}, an
open-source platform that includes the HIP programming model. HIP source code is near-syntactically identical
to CUDA and can be retargeted to NVIDIA via a thin translation
header. However, while HIP simplifies porting CUDA code,
achieving optimal performance on AMD hardware often requires
architecture-specific optimizations that go beyond simple API
translation.

The compression-tool asymmetry between the two ecosystems is
documented in a recent characterization study on AMD GPU
architectures~\cite{anon_iccsa2026}: while NVIDIA has \nvcomp{}
as a production-ready library, the only AMD-side counterpart is
\hipcomp{}, a community-driven port started in 2022 that covers
only LZ4 and Snappy, applies direct source translation, and is
not maintained as a stable API. Scientists working on
AMD-equipped systems including El Capitan, Frontier, and others
cannot directly benefit from the mature CUDA compression
ecosystem and must either port code themselves, use suboptimal
tools, or forgo GPU-accelerated compression entirely. This
situation motivates the present work.

\subsection{Experimental Setup}
\label{sec:bg-setup}

All measurements in this paper use the same Singularity container
(\texttt{arcto\_gfx<arch>.sif}) built per AMD architecture from a
common Ubuntu~22.04 + ROCm~7.0.1 base. \arcto{} is built from
commit \texttt{f8b1c5f} of the \texttt{feature/scaling-instrumentation}
branch. The benchmark binaries time each phase of the
host-to-device path (allocation, pageable-to-pinned memcpy, bulk
H2D, compress kernel, decompress kernel) via host-side
\texttt{std::chrono} and GPU-side \texttt{hipEvent}, exposed by a
\texttt{--report\_phases} flag.

\paragraph{Hardware platforms}
The MI50, MI210, and MI300X experiments use a publicly accessible
HPC testbed at the Luxembourg site (vianden cluster), which
provides exclusive access to single-tenant compute nodes. The
RX~7900~XT experiments run on a workstation-class host at a host
institution research laboratory with a Threadripper PRO 5965WX
processor, 256~GB DDR4, and PCIe~Gen4 x16 host interconnect. Both
platforms run the same container image, differing only in the
architecture-specific build target.

\paragraph{Workloads}
The reference workload is the TTI (Tilted Transversely Isotropic)
seismic wavefield produced by the Fletcher forward-modelling
benchmark~\cite{fletcher2009tti}, the same workload used in the
prior cross-AMD characterization on which this paper
builds~\cite{anon_iccsa2026}. The dataset is
extracted from the middle of a $448^3$ grid simulation with 201
timesteps (to avoid the near-zero values dominant in
initialization timesteps that would inflate compression ratios)
and truncated to six canonical sizes covering the scaling range
of interest: \textbf{10~MB, 100~MB, 1~GB, 4~GB, 8~GB, 16~GB}. Not
every (GPU, size) combination is viable: the smaller-VRAM parts
cannot accommodate the largest inputs because the device input
plus output buffers and the compression scratch space together
exceed the on-package memory. Table~\ref{tab:viability} reports
the (GPU, size) combinations evaluated.

\begin{table}[t]
\centering
\caption{(GPU, size) combinations evaluated in this paper.
\checkmark{} = evaluated; ${\sim}$ = tight (input + output + scratch
near VRAM ceiling, evaluated with caution); --- = exceeds VRAM.
The benchmark allocates a peak of approximately $2 \times$ the input
size on the device (input buffer plus worst-case compressed-output
buffer, both held simultaneously during the compress kernel), so
the maximum input size that fits on a part of VRAM capacity $V$ is
about $V/2$.}
\label{tab:viability}
\footnotesize
\begin{tabular}{l c c c c c c}
\toprule
\textbf{GPU} & \textbf{10 MB} & \textbf{100 MB} & \textbf{1 GB} & \textbf{4 GB} & \textbf{8 GB} & \textbf{16 GB} \\
\midrule
MI50       (32 GB) & \checkmark & \checkmark & \checkmark & \checkmark & ${\sim}$ & --- \\
MI210      (64 GB) & \checkmark & \checkmark & \checkmark & \checkmark & \checkmark & ${\sim}$ \\
MI300X     (192 GB)& \checkmark & \checkmark & \checkmark & \checkmark & \checkmark & \checkmark \\
RX 7900 XT (20 GB) & \checkmark & \checkmark & \checkmark & \checkmark & --- & --- \\
\bottomrule
\end{tabular}
\end{table}

The RX~7900~XT was empirically confirmed to fail at 8~GB with
\texttt{hipErrorOutOfMemory}: the peak simultaneous device
allocation (input buffer at 8~GB plus the worst-case compressed
output buffer at $\approx 8.5$~GB plus the per-chunk scratch)
exceeds the 20~GB VRAM ceiling. The same physical constraint
applies to MI50 at 16~GB. The cross-architecture coverage of the
evaluation is therefore: small inputs validated on all four parts;
large inputs validated on the parts whose VRAM accommodates them.

\paragraph{Metrics and protocol}
Each (algorithm, size, mode) configuration runs in a separate
process invocation with 1 warmup iteration and 3--5 timed
iterations. The benchmark reports per-phase times in milliseconds
(allocation, pageable-to-pinned memcpy, bulk H2D, kernel),
end-to-end time as the sum of phases, and compression ratio.
Cross-mode comparisons use end-to-end speedup of the proposed mode
over the scattered-pageable baseline. The full schema is
documented in the benchmark CSV header and reproduced in the
public repository\footnote{
\TODO{public repo URL}}.


\section{Background and Motivation}
\label{sec:background}

\subsection{GPU-based Data Compression}

Data compression algorithms reduce storage requirements by identifying and eliminating redundancy in data. In the context of scientific computing, lossless algorithms are often preferred as they guarantee exact reconstruction of the original data. Among lossless algorithms, LZ4~\cite{lz4} has gained widespread adoption due to its favorable balance between compression speed and ratio. LZ4 operates by identifying repeated byte sequences within a sliding window and replacing them with compact back-references, a process that can be efficiently parallelized across thousands of GPU threads.

Beyond LZ4, other algorithms offer different trade-offs. Snappy~\cite{snappy}, developed by Google, prioritizes speed over compression ratio and is well-suited for scenarios where minimal latency is critical. Cascaded compression applies multiple encoding stages---including Run-Length Encoding (RLE) and delta encoding---in sequence, achieving higher compression ratios for data with specific patterns at the cost of additional computational overhead.

NVIDIA's nvCOMP library~\cite{nvcomp} has established itself as the reference implementation for GPU-accelerated compression. First released in 2020, nvCOMP provides highly optimized CUDA implementations of LZ4, Snappy, Cascaded, and several other algorithms including zstd, Deflate, and ANS. Through careful optimization of memory access patterns, warp-level primitives, and architecture-specific tuning, nvCOMP achieves throughputs that can saturate the memory bandwidth of high-end GPUs. Version 2.2, which forms the basis of our work, introduced significant performance improvements and remains widely deployed.

\subsection{The AMD GPU Landscape}

AMD has made substantial inroads in high-performance computing, with its GPUs now powering some of the world's fastest supercomputers. This success is built on two distinct GPU architectures optimized for different markets:

\textbf{CDNA (Compute DNA)} targets datacenter and HPC workloads. The architecture emphasizes compute throughput, high-bandwidth memory (HBM), and features specifically designed for scientific computing such as matrix cores for AI workloads. The evolution from CDNA (MI50/MI100) through CDNA~2 (MI210/MI250X) to CDNA~3 (MI300X) has brought substantial improvements in both raw performance and memory bandwidth. The MI300X, in particular, features 304 compute units and 192~GB of HBM3 memory with over 5~TB/s of bandwidth, making it well-suited for memory-intensive workloads like data compression.

\textbf{RDNA (Radeon DNA)} targets gaming and consumer applications but has found use in cost-sensitive HPC deployments. RDNA~3, exemplified by the RX~7900~XT, introduces a chiplet-based design with separate compute and memory dies, achieving competitive compute performance at lower power consumption. While not designed for datacenter use, RDNA GPUs offer an interesting comparison point for understanding how architectural choices affect compression performance.

Table~\ref{tab:gpu_specs} summarizes the key specifications of the GPUs evaluated in this work.

\begin{table}[t]
\centering
\caption{AMD GPU specifications evaluated in this study.}
\label{tab:gpu_specs}
\begin{tabular}{llrrrr}
\toprule
\textbf{GPU} & \textbf{Architecture} & \textbf{CUs} & \textbf{Memory} & \textbf{BW (GB/s)} & \textbf{TDP (W)} \\
\midrule
MI50 & CDNA (Vega 20) & 60 & 16 GB HBM2 & 1,024 & 300 \\
MI210 & CDNA 2 & 104 & 64 GB HBM2e & 1,638 & 300 \\
MI300X & CDNA 3 & 304 & 192 GB HBM3 & 5,300 & 750 \\
RX 7900 XT & RDNA 3 & 84 & 20 GB GDDR6 & 800 & 315 \\
\bottomrule
\end{tabular}
\end{table}

AMD's software ecosystem centers on ROCm (Radeon Open Compute), an open-source platform that provides tools for GPU computing. The HIP (Heterogeneous-compute Interface for Portability) programming model enables developers to write code that can target both AMD and NVIDIA GPUs with minimal modifications, facilitating portability. However, while HIP simplifies porting CUDA code, achieving optimal performance on AMD hardware often requires architecture-specific optimizations that go beyond simple API translation.

\subsection{The Tool Gap}
\label{sec:gap}

Despite AMD's growing presence in HPC, a significant gap exists in the availability of optimized compression tools. Table~\ref{tab:tool_gap} contrasts the state of the NVIDIA and AMD ecosystems for GPU-accelerated compression.

\begin{table}[t]
\centering
\caption{Comparison of GPU compression tool ecosystems.}
\label{tab:tool_gap}
\begin{tabular}{lll}
\toprule
\textbf{Aspect} & \textbf{NVIDIA (nvCOMP)} & \textbf{AMD} \\
\midrule
Primary library & nvCOMP (official) & hipCOMP (community) \\
First release & 2020 & 2022 \\
Maturity & Production-ready & Experimental \\
Algorithms & LZ4, Snappy, Cascaded, & LZ4, Snappy \\
 & zstd, Deflate, ANS, etc. & (limited selection) \\
Optimization level & Highly optimized & Basic ports \\
Documentation & Extensive & Sparse \\
\bottomrule
\end{tabular}
\end{table}

This asymmetry creates practical challenges. Scientists working on AMD-equipped systems---including users of Frontier, LUMI, and other major supercomputers---cannot simply use nvCOMP and must either port code themselves, use suboptimal tools, or forgo GPU-accelerated compression entirely. This situation motivates our work: by characterizing the current state of compression on AMD GPUs and providing an initial implementation, we aim to bridge this gap and enable the broader HPC community to benefit from GPU-accelerated compression regardless of hardware vendor.




%\input{sections/3-implementation}


\section{Porting nvCOMP to AMD GPUs}
\label{sec:implementation}

This section describes our approach to bringing nvCOMP's compression algorithms to AMD GPUs. Our goal was to create a functional implementation that enables performance characterization, prioritizing correctness and compatibility over AMD-specific optimizations.

\subsection{Porting Approach}

We based our implementation on nvCOMP version 2.2, following a methodology similar to the hipCOMP-core project~\cite{hipcomp-core}. The porting process involved three main phases:

\textbf{Phase 1: Automatic Translation.} We used AMD's \texttt{hipify-perl} tool to perform the initial translation of CUDA source code to HIP. This tool automatically converts CUDA API calls to their HIP equivalents (e.g., \texttt{cudaMalloc} $\rightarrow$ \texttt{hipMalloc}, \texttt{cudaMemcpy} $\rightarrow$ \texttt{hipMemcpy}) and adjusts kernel launch syntax. The hipify tool successfully translated the majority of the codebase with minimal manual intervention required for the core compression kernels.

\textbf{Phase 2: Build System Configuration.} The most challenging aspect of the porting effort was configuring the build environment rather than modifying the source code itself. We encountered several issues with CMake configuration, library dependencies, and compiler flags that required careful adjustment. Key modifications included:

\begin{itemize}
    \item Updating CMake files to detect and use HIP compilers (\texttt{hipcc}) instead of CUDA compilers (\texttt{nvcc})
    \item Adjusting library linking to use ROCm equivalents of CUDA runtime libraries
    \item Configuring appropriate compiler flags for AMD GPU architectures (e.g., \texttt{-{}-offload-arch=gfx90a} for MI210, \texttt{-{}-offload-arch=gfx942} for MI300X)
    \item Resolving header file conflicts between CUDA and HIP include paths
\end{itemize}

\textbf{Phase 3: Validation and Testing.} After successful compilation, we validated the implementation by comparing compression outputs against reference data, ensuring that compressed data could be correctly decompressed and that original data was exactly recovered across all test cases.

\subsection{Implementation Environment}

To ensure reproducibility, we containerized our development and execution environment using Singularity~\cite{singularity}. The base environment consists of:

\begin{itemize}
    \item \textbf{ROCm Version}: 7.0.1, installed from the official AMD package
    \item \textbf{Installation}: \texttt{amdgpu-install -{}-usecase=rocm,hip -{}-no-dkms}
    \item \textbf{Container}: Singularity image for portable deployment
\end{itemize}

This containerized approach allows other researchers to reproduce our results without extensive environment configuration, addressing one of the practical barriers to AMD GPU software development.

\subsection{Benchmark Suite}

In addition to porting the compression library itself, we implemented a comprehensive benchmark suite that compiles alongside the main library. The benchmark suite generates synthetic datasets (zeros, binary, random) of configurable sizes, loads scientific datasets (TTI wavefield data), measures throughput with GPU-side timing, validates correctness, and reports statistical metrics across multiple runs. The benchmark structure follows nvCOMP's original design, facilitating methodological consistency with published results.

\subsection{Algorithms Implemented}

Our implementation includes three lossless compression algorithms:

\textbf{LZ4} is a byte-oriented algorithm that achieves high throughput by using a hash-based approach to identify repeated sequences within a sliding window. The GPU implementation parallelizes compression by dividing input data into independent chunks processed concurrently.

\textbf{Snappy}, developed by Google, prioritizes speed over compression ratio using encoding choices that reduce computational overhead at the cost of slightly lower compression ratios.

\textbf{Cascaded} applies multiple encoding stages---specifically Run-Length Encoding (RLE) followed by delta encoding---which is particularly effective for data with local patterns, common in scientific datasets.

\subsection{Current Limitations}

Our implementation is an initial port that prioritizes correctness over optimization:

\begin{itemize}
    \item \textbf{No wavefront-specific tuning}: The code retains CUDA's 32-thread warp assumptions rather than optimizing for AMD's 64-thread wavefronts.
    \item \textbf{No memory hierarchy optimization}: Memory access patterns have not been tuned for AMD's HBM or cache architecture.
    \item \textbf{Limited algorithm coverage}: We focused on three algorithms; nvCOMP supports additional algorithms not yet ported.
\end{itemize}

These limitations are intentional for this characterization study. Future work will address AMD-specific optimizations and extend algorithm coverage.



%\input{sections/4-method}


\section{Experimental Methodology}
\label{sec:methodology}

\subsection{Hardware Platforms}

We evaluated compression performance on four AMD GPUs as detailed in Table~\ref{tab:gpu_specs}, spanning three architectural generations. The MI50 serves as our baseline, representing first-generation CDNA. The MI210 represents CDNA~2, widely deployed in systems like LUMI. The MI300X is AMD's current flagship based on CDNA~3. The RX~7900~XT provides a consumer-grade RDNA~3 comparison point.

All experiments were conducted using ROCm 7.0.1 within Singularity containers to ensure reproducibility.

\subsection{Benchmark Datasets}

We designed our benchmark suite to evaluate compression performance across a range of data characteristics:

\textbf{Synthetic datasets} provide controlled evaluation:
\begin{itemize}
    \item \textbf{Zeros}: All-zero bytes representing highly compressible data
    \item \textbf{Binary}: Random 0/1 patterns with moderate compressibility
    \item \textbf{Random}: Uniformly random bytes establishing an incompressible baseline
\end{itemize}

\textbf{TTI (Tilted Transversely Isotropic)} data represents wavefield snapshots from seismic wave propagation simulations. This dataset is particularly relevant because seismic applications generate massive data volumes, and floating-point scientific data exhibits different compressibility characteristics than synthetic benchmarks.

Each dataset was evaluated at four sizes: \textbf{Small} (10~MB), \textbf{Medium} (100~MB), \textbf{Large} (1~GB), and \textbf{XLarge} (4~GB).

\subsection{Metrics and Procedure}

We report compression and decompression throughput (GB/s), compression ratio (original/compressed size), and speedup relative to the MI50 baseline. For each configuration, we performed 10 timed runs after a warm-up phase and report mean values. Timing used GPU-side events to exclude host-device transfer overhead.



%\input{sections/5-results}



\section{Results: Cross-Generation Analysis}
\label{sec:results}

This section presents our experimental results, analyzing compression performance across AMD GPU generations. We focus on LZ4 as it offers the best balance of throughput and compression ratio for scientific data.

\subsection{Compression Throughput}

Table~\ref{tab:comp_lz4} presents LZ4 compression throughput for the TTI dataset.

\begin{table}[t]
\centering
\caption{LZ4 compression throughput for TTI data (GB/s). Higher is better.}
\label{tab:comp_lz4}
\begin{tabular}{lrrrr}
\toprule
\textbf{GPU} & \textbf{Small} & \textbf{Medium} & \textbf{Large} & \textbf{XLarge} \\
 & (10 MB) & (100 MB) & (1 GB) & (4 GB) \\
\midrule
MI50 (baseline) & 16.0 & 4.7 & 17.9 & 18.7 \\
MI210 & 19.8 & 6.0 & 19.7 & 26.0 \\
MI300X & 34.9 & 8.6 & 24.0 & 40.6 \\
RX 7900 XT & 26.2 & 14.3 & 46.3 & 73.0 \\
\bottomrule
\end{tabular}
\end{table}

Several patterns emerge:

\textbf{Generational improvement within CDNA.} The MI300X achieves 1.3--2.2$\times$ higher compression throughput than the MI50. However, this improvement is modest compared to the 5$\times$ increase in memory bandwidth, suggesting that our unoptimized port does not fully exploit the MI300X's capabilities.

\textbf{RDNA~3 outperforms CDNA~3 for compression.} Surprisingly, the consumer-grade RX~7900~XT achieves higher compression throughput than the datacenter MI300X for medium to xlarge datasets. For XLarge TTI data, the RX~7900~XT reaches 73.0~GB/s compared to 40.6~GB/s---a 1.8$\times$ advantage. This unexpected result suggests that RDNA~3's architecture may be more efficient for certain compression workloads, possibly due to differences in cache behavior or instruction scheduling.

\subsection{Decompression Throughput}

Table~\ref{tab:decomp_lz4} presents decompression results, revealing a markedly different pattern.

\begin{table}[t]
\centering
\caption{LZ4 decompression throughput for TTI data (GB/s). Higher is better.}
\label{tab:decomp_lz4}
\begin{tabular}{lrrrr}
\toprule
\textbf{GPU} & \textbf{Small} & \textbf{Medium} & \textbf{Large} & \textbf{XLarge} \\
 & (10 MB) & (100 MB) & (1 GB) & (4 GB) \\
\midrule
MI50 (baseline) & 53.1 & 25.9 & 26.4 & 27.4 \\
MI210 & 73.1 & 65.4 & 76.6 & 78.8 \\
MI300X & 97.6 & 160.1 & 247.1 & 308.6 \\
RX 7900 XT & 66.0 & 65.9 & 60.1 & 84.5 \\
\bottomrule
\end{tabular}
\end{table}

\textbf{MI300X dominates decompression.} Unlike compression, the MI300X dramatically outperforms all other GPUs. At XLarge size, it achieves 308.6~GB/s---\textbf{11.3$\times$ faster} than the MI50 baseline and 3.7$\times$ faster than the RX~7900~XT. This suggests that decompression better exploits the MI300X's high HBM3 bandwidth.

\textbf{Asymmetric scaling.} While compression speedups were modest (up to 2.2$\times$), decompression speedups reach 11$\times$. This asymmetry has practical implications: workloads that read compressed data more often than they write it will benefit disproportionately from newer CDNA generations.

\subsection{Speedup Summary}

Table~\ref{tab:speedup} summarizes the performance improvements relative to the MI50 baseline.

\begin{table}[t]
\centering
\caption{Speedup vs. MI50 baseline (LZ4, TTI data, XLarge size).}
\label{tab:speedup}
\begin{tabular}{lrr}
\toprule
\textbf{GPU} & \textbf{Compression} & \textbf{Decompression} \\
\midrule
MI210 & 1.4$\times$ & 2.9$\times$ \\
MI300X & 2.2$\times$ & 11.3$\times$ \\
RX 7900 XT & 3.9$\times$ & 3.1$\times$ \\
\bottomrule
\end{tabular}
\end{table}

Key observations:
\begin{itemize}
    \item \textbf{Compression improvements are modest}: Even the best case (RX~7900~XT) achieves only 3.9$\times$ speedup.
    \item \textbf{Decompression scales dramatically}: MI300X achieves 11.3$\times$ speedup.
    \item \textbf{Optimal GPU depends on workload}: RX~7900~XT excels at compression; MI300X excels at decompression.
\end{itemize}

\subsection{Impact of Data Characteristics}

Table~\ref{tab:data_impact} compares throughput across data types on the MI300X.

\begin{table}[t]
\centering
\caption{MI300X compression throughput by data type (GB/s, XLarge).}
\label{tab:data_impact}
\begin{tabular}{lrrr}
\toprule
\textbf{Data Type} & \textbf{LZ4} & \textbf{Snappy} & \textbf{Cascaded} \\
\midrule
Zeros & 844.7 & 72.7 & 402.0 \\
TTI (scientific) & 40.6 & 63.4 & 392.4 \\
\bottomrule
\end{tabular}
\end{table}

LZ4 is highly sensitive to data compressibility: it achieves 844.7~GB/s for zeros but only 40.6~GB/s for TTI data (20$\times$ difference). Cascaded shows more stable performance across data types, making it more predictable for mixed workloads.




%\input{sections/6-discussion}



\section{Discussion}
\label{sec:discussion}

\subsection{Quantifying the Gap}

Our results provide the first comprehensive characterization of compression performance across AMD GPU generations. The performance gap relative to state-of-the-art NVIDIA results is significant but not insurmountable. Published nvCOMP benchmarks report LZ4 throughputs exceeding 200~GB/s on high-end NVIDIA GPUs. Our MI300X results of 40--308~GB/s (compression--decompression) suggest that with optimization, competitive performance may be achievable.

Key factors contributing to the current gap include implementation maturity (nvCOMP has years of optimization), architectural adaptation (our code uses 32-thread assumptions instead of AMD's 64-thread wavefronts), and memory access patterns (not tuned for AMD's cache hierarchy).

\subsection{CDNA vs. RDNA: Unexpected Results}

One of our most surprising findings is that the consumer RDNA~3 GPU outperforms the datacenter CDNA~3 GPU for compression workloads. Possible explanations include RDNA~3's Infinity Cache providing better hit rates for compression's memory access patterns, higher clock speeds on consumer parts benefiting latency-sensitive operations, and RDNA's gaming-optimized architecture handling compression's irregular control flow more efficiently.

These findings suggest that hardware selection for compression workloads should not assume datacenter GPUs are always superior.

\subsection{Practical Implications}

Despite the performance gap, GPU compression on AMD platforms provides value in several scenarios:

\begin{itemize}
    \item \textbf{I/O-bound workloads}: When storage bandwidth is the bottleneck, even moderate compression throughput reduces overall time.
    \item \textbf{Storage cost reduction}: Compression ratios directly reduce storage requirements regardless of throughput.
    \item \textbf{Checkpoint/restart}: Compressing checkpoints reduces I/O time and storage consumption.
\end{itemize}

\subsection{Future Directions}

Our characterization reveals clear optimization opportunities: wavefront-aware kernels for AMD's 64-thread execution model, memory access optimization for AMD's cache hierarchy, and adaptive algorithm selection based on data characteristics. Future work will address these optimizations and extend to lossy compression algorithms such as ZFP for scientific data.



%\input{sections/7-rw}



\section{Related Work}
\label{sec:related}

\textbf{GPU-based compression libraries.} NVIDIA's nvCOMP~\cite{nvcomp} is the most widely used GPU compression library, supporting numerous algorithms with extensive optimization. On the AMD side, hipCOMP-core~\cite{hipcomp-core} provides HIP ports of some algorithms but with limited optimization effort.

\textbf{Scientific data compression.} SZ~\cite{sz} and ZFP~\cite{zfp} provide error-bounded lossy compression for scientific floating-point data. cuSZ~\cite{cusz} brings SZ to NVIDIA GPUs with significant performance improvements. These approaches trade exact reconstruction for higher compression ratios, making them suitable for applications tolerating controlled error.

\textbf{Compression in HPC workflows.} VeloC~\cite{veloc} integrates compression into checkpoint/restart systems. ADIOS~\cite{adios} provides I/O middleware with optional compression. These systems typically rely on nvCOMP for GPU acceleration, highlighting the need for AMD alternatives.

\textbf{AMD GPU performance studies.} Recent work has characterized AMD GPU performance for deep learning and scientific computing, but to our knowledge, no prior study has systematically evaluated compression performance across AMD GPU generations.



%\input{sections/8-conclusion}


\section{Conclusion}
\label{sec:conclusion}

This paper presented the first systematic characterization of GPU-based compression across AMD architectures. Our main findings are:

\begin{enumerate}
    \item \textbf{The tool gap is real but addressable.} AMD platforms lack optimized compression libraries comparable to nvCOMP, but functional implementations are achievable through porting.
    
    \item \textbf{Performance scales with GPU generation.} The MI300X achieves up to 11$\times$ higher decompression throughput than the MI50, demonstrating substantial progress.
    
    \item \textbf{Architecture matters.} CDNA and RDNA show complementary strengths: RDNA~3 excels at compression while CDNA~3 dominates decompression.
    
    \item \textbf{Optimization opportunities abound.} Our unoptimized port leaves significant performance untapped. AMD-specific optimizations could substantially close the gap with nvCOMP.
\end{enumerate}

Future work will focus on AMD-specific optimizations, integration of lossy compression (particularly ZFP), and development of a portable library efficiently targeting both AMD and NVIDIA platforms.





\section*{Acknowledgments}

This study was partially supported by Petrobras grant n.º~2020/00182-5, and by the Coordenação de Aperfeiçoamento de Pessoal de Nível Superior – Brasil (CAPES) – Finance Code 001.

% TODO: Adicionar outros financiamentos/agradecimentos



% \begin{table}
% \caption{Table captions should be placed above the
% tables.}\label{tab1}
% \begin{tabular}{|l|l|l|}
% \hline
% Heading level &  Example & Font size and style\\
% \hline
% Title (centered) &  {\Large\bfseries Lecture Notes} & 14 point, bold\\
% 1st-level heading &  {\large\bfseries 1 Introduction} & 12 point, bold\\
% 2nd-level heading & {\bfseries 2.1 Printing Area} & 10 point, bold\\
% 3rd-level heading & {\bfseries Run-in Heading in Bold.} Text follows & 10 point, bold\\
% 4th-level heading & {\itshape Lowest Level Heading.} Text follows & 10 point, italic\\
% \hline
% \end{tabular}
% \end{table}


% \begin{figure}
% \includegraphics[width=\textwidth]{fig1.eps}
% \caption{A figure caption is always placed below the illustration.
% Please note that short captions are centered, while long ones are
% justified by the macro package automatically.} \label{fig1}
% \end{figure}

% For citations of references, we prefer the use of square brackets
% and consecutive numbers. Citations using labels or the author/year
% convention are also acceptable. The following bibliography provides
% a sample reference list with entries for journal
% articles~\cite{ref_article1}, an LNCS chapter~\cite{ref_lncs1}, a
% book~\cite{ref_book1}, proceedings without editors~\cite{ref_proc1},
% and a homepage~\cite{ref_url1}. Multiple citations are grouped
% \cite{ref_article1,ref_lncs1,ref_book1},
% \cite{ref_article1,ref_book1,ref_proc1,ref_url1}.

%
% ---- Bibliography ----
%
% BibTeX users should specify bibliography style 'splncs04'.
% References will then be sorted and formatted in the correct style.
%

\bibliographystyle{splncs04}
\bibliography{bib/iccsa}

\end{document}
